\section{System analysis}\label{sec:analysis}

This chapter presents a systematic analysis of the problem domain, following established system analysis principles.

\subsection{System representation}

The system under study can be represented as a \enquote{black box} with clearly defined inputs and outputs.
The inputs include raw data, configuration parameters, and user requirements.
The outputs include processed results, performance metrics, and generated reports.

% EXAMPLE: TikZ black-box diagram with labeled I/O
\begin{figure}[htbp]
  \centering
  \begin{tikzpicture}[
      alt={Black-box diagram: raw data, configuration and user requirements flow into the System box; processed results and performance metrics flow out.},
      box/.style={rectangle, draw, thick, minimum width=4cm, minimum height=2.5cm, align=center, fill=gray!10},
      io/.style={->, >=Stealth, thick},
      lbl/.style={font=\small}
    ]
    \node[box] (sys) {\textbf{System}};

    % Inputs
    \node[lbl, left=2cm of sys, yshift=0.7cm] (in1) {Raw data};
    \node[lbl, left=2cm of sys] (in2) {Configuration};
    \node[lbl, left=2cm of sys, yshift=-0.7cm] (in3) {User requirements};

    % Outputs
    \node[lbl, right=2cm of sys, yshift=0.5cm] (out1) {Processed results};
    \node[lbl, right=2cm of sys, yshift=-0.5cm] (out2) {Performance metrics};

    \draw[io] (in1) -- (sys.west |- in1);
    \draw[io] (in2) -- (sys.west |- in2);
    \draw[io] (in3) -- (sys.west |- in3);
    \draw[io] (sys.east |- out1) -- (out1);
    \draw[io] (sys.east |- out2) -- (out2);
  \end{tikzpicture}
  \caption{System \enquote{black box} representation}\label{fig:blackbox}
\end{figure}

\subsection{System functions}

The system performs the following primary functions:
\begin{itemize}
  \item data ingestion and preprocessing;
  \item analysis and transformation;
  \item result generation and visualization;
  \item logging and monitoring;
  \item reporting and export.
\end{itemize}

\subsection{System decomposition}

The system is decomposed into four subsystems:
\begin{itemize}
  \item data processing subsystem;
  \item analysis engine subsystem;
  \item user interface subsystem;
  \item monitoring and logging subsystem.
\end{itemize}

\begin{figure}[htbp]
  \centering
  \begin{tikzpicture}[
      alt={Tree diagram: System splits into Data Processing, Analysis Engine and User Interface; Analysis Engine feeds Monitoring and Logging.},
      box/.style={rectangle, draw, minimum width=3cm, minimum height=0.8cm, align=center},
      >=Stealth
    ]
    \node[box] (sys) {System};
    \node[box, below left=1cm and 0.5cm of sys] (dp) {Data Processing};
    \node[box, below=1cm of sys] (ae) {Analysis Engine};
    \node[box, below right=1cm and 0.5cm of sys] (ui) {User Interface};
    \node[box, below=1cm of ae] (ml) {Monitoring \& Logging};
    \draw[->] (sys) -- (dp);
    \draw[->] (sys) -- (ae);
    \draw[->] (sys) -- (ui);
    \draw[->] (ae) -- (ml);
  \end{tikzpicture}
  \caption{System decomposition diagram}\label{fig:decomposition}
\end{figure}

\subsection{Hierarchical structure}

The hierarchical organization follows the principle of separation of concerns.
At the top level, the orchestrator manages the overall workflow.
The middle layer contains domain-specific processing modules.
The bottom layer handles infrastructure concerns such as data storage and network communication.

\begin{figure}[htbp]
  \centering
  \begin{tikzpicture}[
      alt={Hierarchy diagram: Orchestrator controls Processing and Reporting; Processing uses Storage and Network.},
      box/.style={rectangle, draw, minimum width=2.5cm, minimum height=0.7cm, align=center},
      >=Stealth
    ]
    \node[box] (orch) {Orchestrator};
    \node[box, below left=0.8cm and 1cm of orch] (proc) {Processing};
    \node[box, below right=0.8cm and 1cm of orch] (report) {Reporting};
    \node[box, below left=0.8cm and 0cm of proc] (store) {Storage};
    \node[box, below right=0.8cm and 0cm of proc] (net) {Network};
    \draw[->] (orch) -- (proc);
    \draw[->] (orch) -- (report);
    \draw[->] (proc) -- (store);
    \draw[->] (proc) -- (net);
  \end{tikzpicture}
  \caption{System hierarchy diagram}\label{fig:hierarchy}
\end{figure}

\subsection{State machine}

% EXAMPLE: TikZ state machine / automaton
\cref{fig:state-machine} shows the state transitions of the data processing pipeline.

\begin{figure}[htbp]
  \centering
  \begin{tikzpicture}[
      alt={State machine: Idle goes to Loading on start, Loading to Processing when data is ready, Processing to Done on completion. Loading or Processing go to Error on failure; Error returns to Idle on retry, Done returns to Idle on reset.},
      state/.style={circle, draw, minimum size=1.9cm, align=center, font=\small},
      accept/.style={state, double},
      every node/.append style={font=\small},
      >=Stealth,
    ]
    \node[state] (idle) {Idle};
    \node[state, right=2cm of idle] (load) {Loading};
    \node[state, right=2cm of load] (proc) {Processing};
    \node[accept, right=2cm of proc] (done) {Done};
    % Error sits below the middle of Loading and Processing.
    \node[state] (err) at ($(load)!0.5!(proc) + (0,-3.2cm)$) {Error};

    \draw[->] (idle) -- node[above] {start} (load);
    \draw[->] (load) -- node[above] {data ready} (proc);
    \draw[->] (proc) -- node[above] {complete} (done);
    \draw[->] (load) -- node[left] {fail} (err);
    \draw[->] (proc) -- node[right] {fail} (err);
    % Return arrows leave and enter at fixed angles, clear of the other states.
    \draw[->] (err) to[out=180, in=-90] node[below left] {retry} (idle);
    \draw[->] (done) to[out=110, in=70] node[above] {reset} (idle);
  \end{tikzpicture}
  \caption{Processing pipeline state machine}\label{fig:state-machine}
\end{figure}

\subsection{Resource allocation}

% EXAMPLE: data-driven plot. The numbers live in figures/data/allocation.csv;
% edit the CSV and rebuild, never the picture.
The share of compute resources allocated to each subsystem is shown in \cref{fig:allocation}.

\begin{figure}[htbp]
  \centering
  \begin{tikzpicture}[alt={Horizontal bar chart of resource allocation:
      Processing 40 percent, Analysis 30 percent, User interface 20 percent,
      Monitoring 10 percent.}]
    \begin{axis}[
        xbar,
        width=0.85\linewidth, height=6cm,
        xmin=0, xmax=50,
        xlabel={Share of resources (\%)},
        symbolic y coords={Monitoring, User interface, Analysis, Processing},
        ytick=data,
        nodes near coords,
        nodes near coords style={font=\small},
        bar width=14pt,
      ]
      \addplot[fill=blue!40] table[col sep=comma, x=share, y=subsystem] {figures/data/allocation.csv};
    \end{axis}
  \end{tikzpicture}
  \caption{Resource allocation across subsystems (sample data)}\label{fig:allocation}
\end{figure}

\subsection{Summary}

The system analysis reveals that a modular architecture with clear boundaries between subsystems provides the best foundation for achieving the research objectives.

